\scanpage{159}
\mspage{V.31}
% archive sheet number: 158

où $\alpha$ \deleted{\unsure{est un}} est donné par
\[
\tag*{(6.2.7)}\alpha={\unsure{\kappa|k}}^*h^*,
\]
tandis que $\beta$ est un hom. de degré 1 :
\[
\tag*{(6.2.8)}\unsure{\beta|\rho}:\mathbb{R}\unsure{\Gamma|\mathscr{T}}(V_{\unsure{\eta|y}},\Lambda_{\unsure{\nu|v}})\longrightarrow\mathbb{R}\Gamma(V_s,\Lambda_\nu)(-1)[-1].
\]
\deleted{\uline{\unsure{Proposition}}} Cet homomorphisme peut s'interpréter, par la proposition suivante, qui en \unsure{même temps|toujours} donne un \deleted{\unsure{avec|autre} \unsure{un mode de}} calcul de $\mathbb{R}\Gamma(V,\Lambda_\nu)$, en termes de la donnée fondamentale \deleted{\unsure{en u|en de}} (6.2.1) :

\uline{Proposition} 6.3 a) On a
\deleted{$\underline{H}^i_{V_s}(\unsure{V_s|V},\Lambda_{\nu V})=0$ pour $i\ne2$, et $\underline{H}^2_{V_s}(\Lambda_{\nu V})\simeq\Lambda_{\nu V_s}(1)$,}
\[
\tag*{(6.3.1)}\underline{H}^i_{V_s}(\Lambda_{\nu V})\simeq\begin{cases}
0&\text{si }i\ne2\\
\deleted{\text{si }i}&\\
\Lambda_{\nu V_s}(-1)&\text{si }i=2,
\end{cases}
\]
en d'autres termes :
\[
\tag*{(6.3.2)}\mathbb{R}\underline\Gamma_{V_s}(\Lambda_{\nu V})\simeq\Lambda_{\nu V_s}(-1)[-2],
\]
d'où un isom
\[
\tag*{(6.3.3)}\mathbb{R}\Gamma_{V_s}(V,\Lambda_\nu)\simeq\mathbb{R}\Gamma(V_s,\Lambda_\nu)(-1)[-2]\inserted{\text{ can.}}
\]
b) Considérons le triangle canonique
\diagram{\[
\tag*{(6.3.4)}
\begin{tikzcd}[ampersand replacement=\&,column sep=small]
\& {\mathbb{R}\Gamma(V_\eta,\Lambda_\nu)}\arrow[dl,"\partial"'] \& \\
{\mathbb{R}\Gamma(V_s,\Lambda_\nu)\unsure{(-1)[-2]|[1][-2]}\underset{\text{par (\unsure{6.3.8|6.3.3})}}\simeq\mathbb{R}\Gamma_{V_s}(\deleted{\unsure{V_s}}V,\Lambda_{\nu V})}\arrow[rr,"i_*"] \& \& {\mathbb{R}\Gamma(V,\Lambda_\nu)}\arrow[ul,"j^*"]
\end{tikzcd}
\]}.
Alors la flèche $\partial$ du dit triangle s'identifie [via l'isomorphisme (6.3.3)] à $\beta$. \deleted{Donc $\mathbb{R}\Gamma(V,\Lambda_\nu)$ \unsure{s'identifie à} \unsure{isomorphe}}
\margin{signe ?}
\deleted{$\mathbb{R}\Gamma(V,\Lambda_\nu)$}

\uline{Corollaire} 6.4 \inserted{À des signes et shifts près,} Le triangle (6.3.4) n'est autre que le triangle distingué \deleted{construit avec} \inserted{des} ``{\unsure{mapping-cylindres}}'' du $\beta$ de (6.2.8).
